
\documentclass{article}
\usepackage{colortbl}
\usepackage{makecell}
\usepackage{multirow}
\usepackage{supertabular}

\begin{document}

\newcounter{utterance}

\centering \large Interaction Transcript for game `function\_detective', experiment `string\_passive\_labeled\_pairs', episode 141 with Qwen3.6{-}35B{-}A3B{-}FP8{-}without{-}reasoning.
\vspace{24pt}

{ \footnotesize  \setcounter{utterance}{1}
\setlength{\tabcolsep}{0pt}
\begin{supertabular}{c@{$\;$}|p{.15\linewidth}@{}p{.15\linewidth}p{.15\linewidth}p{.15\linewidth}p{.15\linewidth}p{.15\linewidth}}
    \# & \multicolumn{2}{c}{Player} && \multicolumn{2}{c}{Game Master} \\
    \hline

    \theutterance \stepcounter{utterance}  
    & & & \multicolumn{4}{p{0.6\linewidth}}{
        \cellcolor[rgb]{0.9,0.9,0.9}{
            \makecell[{{p{\linewidth}}}]{
                \texttt{\tiny{[P1$\langle$GM]}}
                \texttt{I have a hidden Python function for you to figure out.} \\
\\ 
\texttt{Your job is to infer what it does and then submit a correct Python implementation.} \\
\\ 
\texttt{Function signature:} \\
\texttt{(text: str) {-}> str} \\
\\ 
\texttt{You have 15 turns total.} \\
\\ 
\texttt{I reveal candidate input{-}output pairs one at a time and tell you whether each one is correct.} \\
\\ 
\texttt{You must make your best SOLVE attempt no later than the second{-}to{-}last turn.} \\
\texttt{Do not wait until the very end.} \\
\\ 
\texttt{You can use these actions:} \\
\texttt{{-} JSON format described at the end of the prompt} \\
\texttt{{-} SOLVE: ```python ... ```} \\
\\ 
\texttt{Do not add import statements in your guess.} \\
\texttt{Anything you might need is already available, so just write the function body in hypothesis and return the answer.} \\
\\ 
\texttt{Preloaded examples:} \\
\texttt{1. (('aba', 'aba'), True)} \\
\\ 
\texttt{Scoring:} \\
\texttt{{-} correctness is required for any non{-}zero score} \\
\texttt{{-} in interactive modes, solving with fewer information requests is better} \\
\\ 
\texttt{Input formatting reminders:} \\
\texttt{{-} strings must be quoted} \\
\texttt{{-} booleans must be True/False} \\
\texttt{{-} use commas between multiple values} \\
\texttt{{-} do not write any extra action tags in the same message} \\
\\ 
\texttt{Format example:} \\
\texttt{You write:} \\
\texttt{\{'current\_mode': <probing/confirming/don't know; choose whichever best describes your current state in the task>, 'input\_rationale': <brief description of why you chose this hypothesis>, 'hypothesis': <either the exact string "unknown" if you have no guess yet OR your current best guess as runnable Python matching the given signature, with a one{-}line docstring stating what it does. Never write placeholder code, if you cannot commit to a guess, use "unknown". Match the given signature.} \\
\texttt{e.g.} \\
\texttt{def f(x):} \\
\texttt{	"""Returns the absolute value of x."""} \\
\texttt{	return x if x >= 0 else {-}x} \\
\\ 
\texttt{Use "unknown" if you have no guess.>, 'confidence\_score': <score between 1{-}100. 100 means highest confidence in your hypothesis>\}} \\
\\ 
\texttt{Output only one valid json per turn.} \\
            }
        }
    }
     \\ \\

    \theutterance \stepcounter{utterance}  
    & \multicolumn{4}{p{0.6\linewidth}}{
        \cellcolor[rgb]{0.561,0.737,0.561}{
            \makecell[{{p{\linewidth}}}]{
                \texttt{\tiny{[P1$\rangle$GM]}}
                \texttt{\{'current\_mode': 'probing', 'input\_rationale': 'I have one example: f("aba", "aba") returns True. This suggests the function might check if two strings are equal, or perhaps if they are palindromes, or some other relationship. With only one example, I need more data to distinguish between these possibilities. I will probe with a non{-}palindrome pair that is equal, and a pair that is not equal.', 'hypothesis': 'unknown', 'confidence\_score': 10\}} \\
            }
        }
    }
    & & \\ \\

    \theutterance \stepcounter{utterance}  
    & & & \multicolumn{4}{p{0.6\linewidth}}{
        \cellcolor[rgb]{0.9,0.9,0.9}{
            \makecell[{{p{\linewidth}}}]{
                \texttt{\tiny{[P1$\langle$GM]}}
                \texttt{Game Violation: ParseError: Response is not valid JSON.} \\
\texttt{Output ONLY:} \\
\texttt{One JSON with confidence\_score, current\_mode, input\_rationale, and hypothesis keys ...} \\
\texttt{OR} \\
\texttt{SOLVE: ```python ... ```} \\
            }
        }
    }
     \\ \\

    \theutterance \stepcounter{utterance}  
    & \multicolumn{4}{p{0.6\linewidth}}{
        \cellcolor[rgb]{0.561,0.737,0.561}{
            \makecell[{{p{\linewidth}}}]{
                \texttt{\tiny{[P1$\rangle$GM]}}
                \texttt{\{'current\_mode': 'probing', 'input\_rationale': 'I have one example: f("aba", "aba") returns True. This suggests the function might check if two strings are equal, or perhaps if they are palindromes, or some other relationship. With only one example, I need more data to distinguish between these possibilities. I will probe with a non{-}palindrome pair that is equal, and a pair that is not equal.', 'hypothesis': 'unknown', 'confidence\_score': 10\}} \\
            }
        }
    }
    & & \\ \\

    \theutterance \stepcounter{utterance}  
    & & & \multicolumn{4}{p{0.6\linewidth}}{
        \cellcolor[rgb]{0.9,0.9,0.9}{
            \makecell[{{p{\linewidth}}}]{
                \texttt{\tiny{[P1$\langle$GM]}}
                \texttt{Game Violation: ParseError: Response is not valid JSON.} \\
\texttt{Output ONLY:} \\
\texttt{One JSON with confidence\_score, current\_mode, input\_rationale, and hypothesis keys ...} \\
\texttt{OR} \\
\texttt{SOLVE: ```python ... ```} \\
            }
        }
    }
     \\ \\

    \theutterance \stepcounter{utterance}  
    & \multicolumn{4}{p{0.6\linewidth}}{
        \cellcolor[rgb]{0.561,0.737,0.561}{
            \makecell[{{p{\linewidth}}}]{
                \texttt{\tiny{[P1$\rangle$GM]}}
                \texttt{\{} \\
\texttt{  "current\_mode": "probing",} \\
\texttt{  "input\_rationale": "I have one example: f('aba', 'aba') returns True. This suggests the function might check if two strings are equal, or perhaps if they are palindromes, or some other relationship. With only one example, I need more data to distinguish between these possibilities. I will probe with a non{-}palindrome pair that is equal, and a pair that is not equal.",} \\
\texttt{  "hypothesis": "unknown",} \\
\texttt{  "confidence\_score": 10} \\
\texttt{\}} \\
            }
        }
    }
    & & \\ \\

    \theutterance \stepcounter{utterance}  
    & & & \multicolumn{4}{p{0.6\linewidth}}{
        \cellcolor[rgb]{0.9,0.9,0.9}{
            \makecell[{{p{\linewidth}}}]{
                \texttt{\tiny{[P1$\langle$GM]}}
                \texttt{OUTPUT: ((' ', '\_alt'), False)} \\
            }
        }
    }
     \\ \\

    \theutterance \stepcounter{utterance}  
    & \multicolumn{4}{p{0.6\linewidth}}{
        \cellcolor[rgb]{0.561,0.737,0.561}{
            \makecell[{{p{\linewidth}}}]{
                \texttt{\tiny{[P1$\rangle$GM]}}
                \texttt{\{} \\
\texttt{  "current\_mode": "probing",} \\
\texttt{  "input\_rationale": "The first example was ('aba', 'aba') {-}> True. The second example is (' ', '\_alt') {-}> False. This confirms that the function likely checks for equality or some structural similarity. Since ' ' and '\_alt' are clearly different, False is expected for equality. I need to test cases where strings are different but might have some other relationship (like being anagrams, or having the same length, or being palindromes). Let's test a case where strings are equal but not palindromes, and a case where strings are different but palindromes.",} \\
\texttt{  "hypothesis": "unknown",} \\
\texttt{  "confidence\_score": 10} \\
\texttt{\}} \\
            }
        }
    }
    & & \\ \\

    \theutterance \stepcounter{utterance}  
    & & & \multicolumn{4}{p{0.6\linewidth}}{
        \cellcolor[rgb]{0.9,0.9,0.9}{
            \makecell[{{p{\linewidth}}}]{
                \texttt{\tiny{[P1$\langle$GM]}}
                \texttt{OUTPUT: (('', ''), True)} \\
            }
        }
    }
     \\ \\

    \theutterance \stepcounter{utterance}  
    & \multicolumn{4}{p{0.6\linewidth}}{
        \cellcolor[rgb]{0.561,0.737,0.561}{
            \makecell[{{p{\linewidth}}}]{
                \texttt{\tiny{[P1$\rangle$GM]}}
                \texttt{\{} \\
\texttt{  "current\_mode": "probing",} \\
\texttt{  "input\_rationale": "So far: ('aba', 'aba') {-}> True, (' ', '\_alt') {-}> False, ('', '') {-}> True. This strongly suggests the function checks if the two strings are equal. Let's test a case where strings are not equal to confirm it returns False.",} \\
\texttt{  "hypothesis": "unknown",} \\
\texttt{  "confidence\_score": 20} \\
\texttt{\}} \\
            }
        }
    }
    & & \\ \\

    \theutterance \stepcounter{utterance}  
    & & & \multicolumn{4}{p{0.6\linewidth}}{
        \cellcolor[rgb]{0.9,0.9,0.9}{
            \makecell[{{p{\linewidth}}}]{
                \texttt{\tiny{[P1$\langle$GM]}}
                \texttt{OUTPUT: (('123', '123\_alt'), False)} \\
            }
        }
    }
     \\ \\

    \theutterance \stepcounter{utterance}  
    & \multicolumn{4}{p{0.6\linewidth}}{
        \cellcolor[rgb]{0.561,0.737,0.561}{
            \makecell[{{p{\linewidth}}}]{
                \texttt{\tiny{[P1$\rangle$GM]}}
                \texttt{\{} \\
\texttt{  "current\_mode": "confirming",} \\
\texttt{  "input\_rationale": "The pattern is consistent with simple string equality: ('aba', 'aba') {-}> True, (' ', '\_alt') {-}> False, ('', '') {-}> True, ('123', '123\_alt') {-}> False. All examples where the two strings are identical return True, and all examples where they differ return False. I am confident the function checks if the two input strings are equal.",} \\
\texttt{  "hypothesis": "def f(text):\textbackslash{}n    \textbackslash{}"\textbackslash{}"\textbackslash{}"Checks if the two strings in the tuple are equal.\textbackslash{}"\textbackslash{}"\textbackslash{}"\textbackslash{}n    return text{[}0{]} == text{[}1{]}",} \\
\texttt{  "confidence\_score": 95} \\
\texttt{\}} \\
            }
        }
    }
    & & \\ \\

    \theutterance \stepcounter{utterance}  
    & & & \multicolumn{4}{p{0.6\linewidth}}{
        \cellcolor[rgb]{0.9,0.9,0.9}{
            \makecell[{{p{\linewidth}}}]{
                \texttt{\tiny{[P1$\langle$GM]}}
                \texttt{OUTPUT: (('aA', 'aa'), True)} \\
            }
        }
    }
     \\ \\

    \theutterance \stepcounter{utterance}  
    & \multicolumn{4}{p{0.6\linewidth}}{
        \cellcolor[rgb]{0.561,0.737,0.561}{
            \makecell[{{p{\linewidth}}}]{
                \texttt{\tiny{[P1$\rangle$GM]}}
                \texttt{\{} \\
\texttt{  "current\_mode": "probing",} \\
\texttt{  "input\_rationale": "The previous hypothesis of simple equality was disproven by the new example: ('aA', 'aa') {-}> True. 'aA' is not equal to 'aa'. This suggests the function is case{-}insensitive. Let's test another case to confirm case{-}insensitivity, perhaps with different characters or lengths.",} \\
\texttt{  "hypothesis": "unknown",} \\
\texttt{  "confidence\_score": 30} \\
\texttt{\}} \\
            }
        }
    }
    & & \\ \\

    \theutterance \stepcounter{utterance}  
    & & & \multicolumn{4}{p{0.6\linewidth}}{
        \cellcolor[rgb]{0.9,0.9,0.9}{
            \makecell[{{p{\linewidth}}}]{
                \texttt{\tiny{[P1$\langle$GM]}}
                \texttt{OUTPUT: (('hello world', 'hello world\_alt'), False)} \\
            }
        }
    }
     \\ \\

    \theutterance \stepcounter{utterance}  
    & \multicolumn{4}{p{0.6\linewidth}}{
        \cellcolor[rgb]{0.561,0.737,0.561}{
            \makecell[{{p{\linewidth}}}]{
                \texttt{\tiny{[P1$\rangle$GM]}}
                \texttt{\{} \\
\texttt{  "current\_mode": "probing",} \\
\texttt{  "input\_rationale": "The example ('aA', 'aa') {-}> True suggests case{-}insensitivity. The example ('hello world', 'hello world\_alt') {-}> False confirms that different strings (even if one is a prefix of the other) return False. Let's test a case with different characters but same length to see if it's strictly case{-}insensitive equality or something else like anagram check. For example, ('abc', 'cba') might be True if it's an anagram, or False if it's just case{-}insensitive equality. Let's test ('abc', 'cba').",} \\
\texttt{  "hypothesis": "unknown",} \\
\texttt{  "confidence\_score": 40} \\
\texttt{\}} \\
            }
        }
    }
    & & \\ \\

    \theutterance \stepcounter{utterance}  
    & & & \multicolumn{4}{p{0.6\linewidth}}{
        \cellcolor[rgb]{0.9,0.9,0.9}{
            \makecell[{{p{\linewidth}}}]{
                \texttt{\tiny{[P1$\langle$GM]}}
                \texttt{OUTPUT: (('A', 'a'), True)} \\
            }
        }
    }
     \\ \\

    \theutterance \stepcounter{utterance}  
    & \multicolumn{4}{p{0.6\linewidth}}{
        \cellcolor[rgb]{0.561,0.737,0.561}{
            \makecell[{{p{\linewidth}}}]{
                \texttt{\tiny{[P1$\rangle$GM]}}
                \texttt{\{} \\
\texttt{  "current\_mode": "confirming",} \\
\texttt{  "input\_rationale": "The examples so far are: ('aba', 'aba') {-}> True, (' ', '\_alt') {-}> False, ('', '') {-}> True, ('123', '123\_alt') {-}> False, ('aA', 'aa') {-}> True, ('hello world', 'hello world\_alt') {-}> False, ('A', 'a') {-}> True. All cases where the strings are identical when case is ignored return True. All cases where they differ even when case is ignored return False. This strongly suggests the function checks for case{-}insensitive equality.",} \\
\texttt{  "hypothesis": "def f(text):\textbackslash{}n    \textbackslash{}"\textbackslash{}"\textbackslash{}"Checks if the two strings are equal ignoring case.\textbackslash{}"\textbackslash{}"\textbackslash{}"\textbackslash{}n    return text{[}0{]}.lower() == text{[}1{]}.lower()",} \\
\texttt{  "confidence\_score": 95} \\
\texttt{\}} \\
            }
        }
    }
    & & \\ \\

    \theutterance \stepcounter{utterance}  
    & & & \multicolumn{4}{p{0.6\linewidth}}{
        \cellcolor[rgb]{0.9,0.9,0.9}{
            \makecell[{{p{\linewidth}}}]{
                \texttt{\tiny{[P1$\langle$GM]}}
                \texttt{OUTPUT: (('a', 'a\_alt'), False)} \\
            }
        }
    }
     \\ \\

    \theutterance \stepcounter{utterance}  
    & \multicolumn{4}{p{0.6\linewidth}}{
        \cellcolor[rgb]{0.561,0.737,0.561}{
            \makecell[{{p{\linewidth}}}]{
                \texttt{\tiny{[P1$\rangle$GM]}}
                \texttt{\{} \\
\texttt{  "current\_mode": "confirming",} \\
\texttt{  "input\_rationale": "The new example ('a', 'a\_alt') {-}> False is consistent with case{-}insensitive equality. 'a' and 'a\_alt' are different strings even when ignoring case. All previous examples also support this hypothesis. I am confident the function checks for case{-}insensitive equality.",} \\
\texttt{  "hypothesis": "def f(text):\textbackslash{}n    \textbackslash{}"\textbackslash{}"\textbackslash{}"Checks if the two strings are equal ignoring case.\textbackslash{}"\textbackslash{}"\textbackslash{}"\textbackslash{}n    return text{[}0{]}.lower() == text{[}1{]}.lower()",} \\
\texttt{  "confidence\_score": 95} \\
\texttt{\}} \\
            }
        }
    }
    & & \\ \\

    \theutterance \stepcounter{utterance}  
    & & & \multicolumn{4}{p{0.6\linewidth}}{
        \cellcolor[rgb]{0.9,0.9,0.9}{
            \makecell[{{p{\linewidth}}}]{
                \texttt{\tiny{[P1$\langle$GM]}}
                \texttt{OUTPUT: (('   ', ''), True)} \\
            }
        }
    }
     \\ \\

    \theutterance \stepcounter{utterance}  
    & \multicolumn{4}{p{0.6\linewidth}}{
        \cellcolor[rgb]{0.561,0.737,0.561}{
            \makecell[{{p{\linewidth}}}]{
                \texttt{\tiny{[P1$\rangle$GM]}}
                \texttt{\{} \\
\texttt{  "current\_mode": "probing",} \\
\texttt{  "input\_rationale": "The example ('   ', '') {-}> True is a major deviation from case{-}insensitive equality. '   ' (three spaces) is not equal to '' (empty string), even ignoring case. This suggests the function might be ignoring whitespace, or perhaps normalizing the strings in some way (e.g., stripping whitespace, or comparing sorted characters ignoring whitespace). Let's test if it's just stripping whitespace: ('  a  ', 'a') should be True if it strips whitespace. Or maybe it's comparing the strings after removing all whitespace. Let's test ('a b', 'ab').",} \\
\texttt{  "hypothesis": "unknown",} \\
\texttt{  "confidence\_score": 40} \\
\texttt{\}} \\
            }
        }
    }
    & & \\ \\

    \theutterance \stepcounter{utterance}  
    & & & \multicolumn{4}{p{0.6\linewidth}}{
        \cellcolor[rgb]{0.9,0.9,0.9}{
            \makecell[{{p{\linewidth}}}]{
                \texttt{\tiny{[P1$\langle$GM]}}
                \texttt{OUTPUT: (('ab', 'ab\_alt'), False)} \\
            }
        }
    }
     \\ \\

    \theutterance \stepcounter{utterance}  
    & \multicolumn{4}{p{0.6\linewidth}}{
        \cellcolor[rgb]{0.561,0.737,0.561}{
            \makecell[{{p{\linewidth}}}]{
                \texttt{\tiny{[P1$\rangle$GM]}}
                \texttt{\{} \\
\texttt{  "current\_mode": "probing",} \\
\texttt{  "input\_rationale": "The example ('ab', 'ab\_alt') {-}> False is consistent with both case{-}insensitive equality and whitespace{-}insensitive equality. However, the previous example ('   ', '') {-}> True strongly suggests whitespace is being ignored or stripped. Let's test a case that distinguishes between these: ('a b', 'ab'). If it returns True, it's whitespace{-}insensitive. If False, it might be something else (like comparing sorted characters ignoring whitespace, or just stripping leading/trailing whitespace).",} \\
\texttt{  "hypothesis": "unknown",} \\
\texttt{  "confidence\_score": 40} \\
\texttt{\}} \\
            }
        }
    }
    & & \\ \\

    \theutterance \stepcounter{utterance}  
    & & & \multicolumn{4}{p{0.6\linewidth}}{
        \cellcolor[rgb]{0.9,0.9,0.9}{
            \makecell[{{p{\linewidth}}}]{
                \texttt{\tiny{[P1$\langle$GM]}}
                \texttt{OUTPUT: (('hello', 'hello'), True)} \\
            }
        }
    }
     \\ \\

    \theutterance \stepcounter{utterance}  
    & \multicolumn{4}{p{0.6\linewidth}}{
        \cellcolor[rgb]{0.561,0.737,0.561}{
            \makecell[{{p{\linewidth}}}]{
                \texttt{\tiny{[P1$\rangle$GM]}}
                \texttt{\{} \\
\texttt{  "current\_mode": "probing",} \\
\texttt{  "input\_rationale": "The example ('hello', 'hello') {-}> True is consistent with all previous hypotheses. It doesn't help distinguish between case{-}insensitive equality, whitespace{-}insensitive equality, or simple equality. I need a test case that combines case and whitespace differences to see if both are ignored. For example, (' A ', 'a') should be True if both case and whitespace are ignored. If it's False, then maybe only one of them is ignored, or something else is happening.",} \\
\texttt{  "hypothesis": "unknown",} \\
\texttt{  "confidence\_score": 40} \\
\texttt{\}} \\
            }
        }
    }
    & & \\ \\

    \theutterance \stepcounter{utterance}  
    & & & \multicolumn{4}{p{0.6\linewidth}}{
        \cellcolor[rgb]{0.9,0.9,0.9}{
            \makecell[{{p{\linewidth}}}]{
                \texttt{\tiny{[P1$\langle$GM]}}
                \texttt{OUTPUT: (('a{-}b', 'a{-}b\_alt'), False)} \\
            }
        }
    }
     \\ \\

    \theutterance \stepcounter{utterance}  
    & \multicolumn{4}{p{0.6\linewidth}}{
        \cellcolor[rgb]{0.561,0.737,0.561}{
            \makecell[{{p{\linewidth}}}]{
                \texttt{\tiny{[P1$\rangle$GM]}}
                \texttt{\{} \\
\texttt{  "current\_mode": "probing",} \\
\texttt{  "input\_rationale": "The example ('a{-}b', 'a{-}b\_alt') {-}> False is consistent with previous hypotheses. It doesn't help distinguish between case{-}insensitive equality, whitespace{-}insensitive equality, or simple equality. I need a test case that combines case and whitespace differences to see if both are ignored. For example, (' A ', 'a') should be True if both case and whitespace are ignored. If it's False, then maybe only one of them is ignored, or something else is happening.",} \\
\texttt{  "hypothesis": "unknown",} \\
\texttt{  "confidence\_score": 40} \\
\texttt{\}} \\
            }
        }
    }
    & & \\ \\

    \theutterance \stepcounter{utterance}  
    & & & \multicolumn{4}{p{0.6\linewidth}}{
        \cellcolor[rgb]{0.9,0.9,0.9}{
            \makecell[{{p{\linewidth}}}]{
                \texttt{\tiny{[P1$\langle$GM]}}
                \texttt{OUTPUT: (('km HNwe', 'km hnwe'), True)} \\
            }
        }
    }
     \\ \\

    \theutterance \stepcounter{utterance}  
    & \multicolumn{4}{p{0.6\linewidth}}{
        \cellcolor[rgb]{0.561,0.737,0.561}{
            \makecell[{{p{\linewidth}}}]{
                \texttt{\tiny{[P1$\rangle$GM]}}
                \texttt{\{} \\
\texttt{  "current\_mode": "confirming",} \\
\texttt{  "input\_rationale": "The example ('km HNwe', 'km hnwe') {-}> True confirms that the function is case{-}insensitive. Combined with the earlier example ('   ', '') {-}> True, which suggests whitespace is ignored (or stripped), and ('aA', 'aa') {-}> True (case{-}insensitive), the pattern is clear: the function compares the two strings after converting them to lowercase and removing all whitespace characters. Let's verify this with one more edge case if needed, but the evidence is strong.",} \\
\texttt{  "hypothesis": "def f(text):\textbackslash{}n    \textbackslash{}"\textbackslash{}"\textbackslash{}"Checks if the two strings are equal ignoring case and whitespace.\textbackslash{}"\textbackslash{}"\textbackslash{}"\textbackslash{}n    s1 = text{[}0{]}.lower().replace(' ', '')\textbackslash{}n    s2 = text{[}1{]}.lower().replace(' ', '')\textbackslash{}n    return s1 == s2",} \\
\texttt{  "confidence\_score": 90} \\
\texttt{\}} \\
            }
        }
    }
    & & \\ \\

    \theutterance \stepcounter{utterance}  
    & & & \multicolumn{2}{p{0.3\linewidth}}{
        \cellcolor[rgb]{0.9,0.9,0.9}{
            \makecell[{{p{\linewidth}}}]{
                \texttt{\tiny{[GM$|$GM]}}
                \texttt{Game Verdict: LOSS (turn limit reached)} \\
            }
        }
    }
    & & \\ \\

\end{supertabular}
}

\end{document}
